\section*{अध्याय \ref{ch:g8:sound-pitch-loudness} --- ध्वनि: चाल, तारत्व, प्रबलता}

\begin{solution}{exo:g8:sound-pitch-loudness:1}
\omterm{def:g8:sound-pitch-loudness:frequency}{आवृत्ति}: हर सेकंड लगने वाले पूरे झूलों की गिनती, \omterm{def:g8:sound-pitch-loudness:frequency}{हर्ट्ज़} (\unit{Hz})
में। मंच वाला A: \qty{440}{Hz}; और उससे एक सप्तक ऊपर:
\qty{880}{Hz} --- यानी दोगुना।
\end{solution}

\begin{solution}{exo:g8:sound-pitch-loudness:2}
\qty{0.02}{s} में $4$ झूले: \qty{200}{Hz}; और $16$ झूले:
\qty{800}{Hz} --- दूसरी ऊँची सुनाई देगी (दरअसल पूरे दो सप्तक ऊपर)।
\end{solution}

\begin{solution}{exo:g8:sound-pitch-loudness:3}
क़रीब \qty{20}{Hz} से \qty{20000}{Hz} तक। नीचे: \omterm{prop:g8:sound-pitch-loudness:range}{अवश्रव्य ध्वनि} ---
हाथियों की दूर तक जाती गड़गड़ाहटें। ऊपर: \omterm{prop:g8:sound-pitch-loudness:range}{पराश्रव्य ध्वनि} --- चमगादड़
की शिकार वाली चटकें।
\end{solution}

\begin{solution}{exo:g8:sound-pitch-loudness:4}
वही सुर, ज़्यादा तेज़: शिखर ऊँचे हो जाते हैं, फ़ासला जस का तस। और
उतनी ही तेज़ आवाज़ में ऊँचा सुर: शिखर पास-पास सिमट आते हैं, ऊँचाई जस
की तस। दो आज़ाद घुंडियाँ।
\end{solution}

\begin{solution}{exo:g8:sound-pitch-loudness:5}
हर एक \omterm{def:g8:sound-pitch-loudness:frequency}{आवृत्ति} बढ़ा देता है: छोटे, कसे हुए और पतले तार हर सेकंड
ज़्यादा बार काँपते हैं --- बड़ी $f$, ऊँचा तारत्व।
\end{solution}

\begin{solution}{exo:g8:sound-pitch-loudness:6}
पत्ते \qty{10}{dB}; यातायात \qty{80}{dB}; अगली क़तार \qty{110}{dB};
दर्द क़रीब \qty{120}{dB}। नुक़सान वाली पट्टी क़रीब \qty{85}{dB} से
खुलती है।
\end{solution}

\begin{solution}{exo:g8:sound-pitch-loudness:7}
उस सीटी का तारत्व \qty{20000}{Hz} से ज़रा ही ऊपर बैठता है ---
मालिक की खिड़की के बाहर, और कुत्ते की खिड़की के भीतर। एक जोड़ी कानों
के लिए ख़राब, दूसरी के लिए चुभने वाली: खिड़कियाँ प्रजाति दर प्रजाति
अलग होती हैं।
\end{solution}

\begin{solution}{exo:g8:sound-pitch-loudness:8}
यह कि वह \omterm{def:g7:motion-average-speed:formula}{चाल} हर तारत्व और हर प्रबलता के लिए एक ही है: अलग-अलग चालें
दूर के समस्वरों को बिखेरकर आरोह बना देतीं और धीमी आवाज़ को तेज़ के
पीछे छोड़ देतीं --- पार्क का कार्यक्रम रोज़ रात दोनों को झुठला देता
है।
\end{solution}

\begin{solution}{exo:g8:sound-pitch-loudness:9}
हर चटक में $50000 \times 0.01 = 500$ झूले। बारीक चित्रकारी को बारीक
क़लम चाहिए: सिर्फ़ बहुत छोटे झूलों वाले \omterm{def:g3:sound-around-us:vibration}{कंपन} --- यानी बहुत बड़ी $f$
--- ही अपनी प्रतिध्वनियों में छोटी-छोटी बारीकियाँ टटोल सकते हैं; जबकि
मोटी, कम \omterm{def:g8:sound-pitch-loudness:frequency}{आवृत्ति} वाली गड़गड़ाहटें किसी पतिंगे के ऊपर से बिना पहचाने ही
बह जाती हैं।
\end{solution}

\begin{solution}{exo:g8:sound-pitch-loudness:10}
दोनों कुंजियाँ खिड़की के भीतर बैठती हैं --- \qty{27}{Hz}
\qty{20}{Hz} वाले फ़र्श से ज़रा ऊपर, और \qty{4200}{Hz} छत से आराम से
नीचे। पर उम्र इस खिड़की को \emph{ऊपर} से घिसती है: मिलाने वाले का अपना
कान पहले ऊँचे सुरों में मंद पड़ता है, इसलिए जाँचने वाला यंत्र सबसे
नीचे वाले सप्तकों से पहले सबसे ऊँचे सप्तकों के लिए बुलाया जाता है।
\end{solution}

\begin{solution}{exo:g8:sound-pitch-loudness:11}
बराबर क़दमों में दस गुना बढ़त छिपी है --- हर कुछ डंडे \omterm{def:g3:sound-around-us:sound}{ध्वनि} की ताक़त
जोड़ते नहीं, गुणा करते हैं। इसलिए सीढ़ी पर ``ज़रा-सा और तेज़'' कान के
लिए \emph{कहीं} ज़्यादा तगड़ा होता है: जो घुंडी की छोटी-सी घुमाई जैसा
पढ़ा जाता है, उसी पर बजट \omterm{ex:g2:measuring-time:clock}{घंटों} से मिनटों में ढह जाता है।
\end{solution}

\begin{solution}{exo:g8:sound-pitch-loudness:12}
पियानो की जाँच: यंत्र तार के बग़ल रखकर उसका पाठ $440$ से मिलाओ। पिछली
क़तार की प्रबलता: ढोल की परख के दौरान वहीं खड़े होकर \omterm{ex:g8:sound-pitch-loudness:decibels}{डेसिबल} का पाठ
पढ़ो। \omterm{def:g2:air-around-us:air}{हवा} वाली मशीन: यंत्र का \omterm{prop:g8:colors-spectra:mixture}{वर्णक्रम} वाला पर्दा नीचे वाले सिरे पर
ताको और \qty{20}{Hz} से नीचे कोई चोटी ढूँढ़ो। सबसे कम भरोसेमंद:
अवश्रव्य वाला फ़ैसला --- फ़ोन के माइक्रोफ़ोन आवाज़ों के लिए बने होते
हैं और खिड़की के फ़र्श के पास बहरे हो जाते हैं, इसलिए वहाँ ख़ामोश पर्दा
कुछ ख़ास साबित नहीं करता।
\end{solution}

\begin{solution}{pb:g8:sound-pitch-loudness:1}
\textbf{1.} पियानो उतरा हुआ है --- उसका हर A हर सेकंड आठ झूले कम
काँप रहा है। घुंडी तारत्व वाली है: यानी मिलाने वाले की चाबी, कोई
प्रवर्धक नहीं।
\textbf{2.} \qty{41}{Hz} खिड़की के \qty{20}{Hz} वाले फ़र्श से ज़रा ही
ऊपर बैठती है: सुनाई देगी, गहरी --- और इतने धीमे झूलों पर छाती का अपना
ढोल भी साथ जवाब देता है: सुनने जितना ही महसूस भी।
\textbf{3.} वही तारत्व (बराबर फ़ासला, \qty{440}{Hz}), और तिगुना \omterm{prop:g8:sound-pitch-loudness:loudness}{आयाम}:
गायक द्विभुज का ही सुर गा रहा है, बस कहीं ज़्यादा तेज़।
\textbf{4.} सिर्फ़ सबसे जवान कान --- \qty{20000}{Hz} के पास वाली
खिड़की की छत बच्चों की मिलकियत है; ज़्यादातर बड़ों की खिड़कियाँ पहले
ही नीचे उतर चुकी हैं, और \qty{22000}{Hz} की तसदीक़ कोई सुनने वाला नहीं
कर सकता।
\textbf{5.} रॉक वाला हिस्सा (\qty{100}{dB}) और समापन (\qty{110}{dB}),
दोनों \qty{85}{dB} वाली नुक़सान की रेखा से ऊपर खड़े हैं: यानी सिर्फ़
मिनटों वाला इलाक़ा; जबकि बिना बिजली वाला हिस्सा पास हो जाता है।
\textbf{6.} ठेपों के पीछे वह क़रीब \qty{80}{dB} बन जाता है ---
वापस नुक़सान की रेखा के नीचे: पूरी शाम का सामना झेलने लायक बन गया,
और इसीलिए समझदार लोग उन्हें हमेशा लगाते हैं।
\textbf{7.} फैलते खोल: दूरी के साथ हर ठसाठस और विशाल खोल पर बँटती जाती
है --- यानी प्रबलता (\omterm{prop:g8:sound-pitch-loudness:loudness}{आयाम}) फीकी पड़ती है। तारत्व जस का तस सवार रहता
है: दूरी सिर्फ़ एक ही घुंडी घुमाती है।
\textbf{8.} धीरे बजाना सिर्फ़ \omterm{prop:g8:sound-pitch-loudness:loudness}{आयाम} वाली घुंडी घुमाता है --- हर सुर की
\omterm{def:g8:sound-pitch-loudness:frequency}{आवृत्ति}, ऊँची हो या नीची, ठीक अपनी जगह टिकी रहती है। (उसका ``फीका''
कम प्रबलता पर ख़ुद कान की तरकीब है, तारों की भौतिकी नहीं।)
\textbf{9.} $680 \div 340 = \qty{2}{s}$ की कौंध-से-धमाके वाली देरी।
और हर सेकंड दो धमक: \qty{2}{Hz} --- फ़र्श से बहुत नीचे: छाती में
महसूस होगी, सुर की तरह कभी सुनाई नहीं देगी।
\textbf{10.} बास के सबसे नीचे वाले झूले: खिड़की के फ़र्श के पास और
उससे नीचे, ठीक वहीं जहाँ आवाज़ों के लिए बना माइक्रोफ़ोन बहरा हो जाता
है --- वह गड़गड़ाहट सचमुच थी; फ़ोन ने उसे कभी पकड़ा ही नहीं।
\textbf{11.} जैसे: ``तारत्व \omterm{def:g8:sound-pitch-loudness:frequency}{आवृत्ति} है --- हर सेकंड के झूलों की
गिनती, और सुर मिलाने वाले का पूरा कारोबार। प्रबलता \omterm{prop:g8:sound-pitch-loudness:loudness}{आयाम} है ---
झूलों का नाप, और नर्स का पूरा बहीखाता। और \omterm{def:g7:motion-average-speed:formula}{चाल} को कोई भी घुंडी नहीं
छूती: आज रात हर सुर ने वह हॉल उसी \qty{340}{m/s} से पार किया।''
\textbf{12.} खिड़की वाली प्रतिज्ञप्ति: यंत्र का सुर \qty{20000}{Hz}
से ऊपर --- यानी हर इंसानी खिड़की के बाहर, और चूहे की खिड़की के भीतर।
एक ही \omterm{def:g3:sound-around-us:vibration}{कंपन}, दो दर्शक-वर्ग, और कोई विरोधाभास नहीं।
\end{solution}
